% Source: https://artofproblemsolving.com/community/c3h2668847_geo_problem
% Stable identity: aops-topic-2668847-q01
Put the centre of the rhombus at the origin: $A=(-4,0)$, $C=(4,0)$, $B=(0,3)$, $D=(0,-3)$. Line $BC$ is $3x+4y=12$.

A circle tangent to lines $BC$ and $CD$ has its centre on a bisector of the angle between them: the line $AC$ or the line $x=4$. A centre $(4,k)$ would have radius $\dfrac{4|k|}{5}$, less than its distance $\sqrt{64+k^2}$ to $A$; impossible. So the centre is $I=(h,0)$, with radius $|h+4|$ (the distance to $A$) equal to the distance $\dfrac{|3h-12|}{5}$ to line $BC$.
\begin{itemize}
\item $-4<h<4$: $5(h+4)=12-3h$, so $h=-1$ and $r=3$. The tangency point on line $BC$ is $(0.8,2.4)$, which lies on segment $BC$ (and symmetrically on $CD$).
\item $h\ge4$: $5(h+4)=3h-12$ gives $h=-16$, a contradiction.
\item $h<-4$: $-5(h+4)=12-3h$ gives $h=-16$, $r=12$, but then the circle touches line $BC$ at $(-8.8,9.6)$, outside the segment.
\end{itemize}
So the radius is $\boxed{3}$.

\begin{asy}
usepackage("amsmath");
size(3.5cm);
defaultpen(fontsize(10pt));
pair A=(-4,0),B=(0,3),C=(4,0),D=(0,-3),O=(-1,0),T1=(0.8,2.4),T2=(0.8,-2.4);
draw(A--B--C--D--cycle,linewidth(0.8));draw(A--C,dashed);draw(B--D,dashed);draw(circle(O,3),linewidth(0.7));
draw(O--T1,dotted);draw(O--T2,dotted);dot(O);dot(T1);dot(T2);
label("$A$",A,W);label("$B$",B,N);label("$C$",C,E);label("$D$",D,S);label("$I$",O,SW,fontsize(8pt));
\end{asy}
